%~Mouliné par MaN_auto v.0.21.8 2019-11-22 17:01:41
\documentclass[AHL,Unicode,longabstracts]{cedram}

\datereceived{2018-02-27}
\daterevised{2018-10-03}
\dateaccepted{2019-05-09}
\dateposted{2020-01-16}

\editor{P. G\'erard}

\newcommand{\BA}{\mathbf{A}}
\newcommand{\BR}{\mathbf{R}}

\newcommand{\R}{\mathbb{R}}
\newcommand{\Com}{\mathbb{C}}
\newcommand{\N}{\mathbb{N}}

%\newcommand{\IM}{\,\mathrm{Im}}
%\newcommand{\Div}{\,\mathrm{div}}

\DeclareMathOperator{\IM}{Im}
\DeclareMathOperator{\Div}{div}
\DeclareMathOperator{\RE}{Re}

\newcommand*{\dd}{\mathrm{d}}
\newcommand*{\ds}{\dd s}
\newcommand*{\dx}{\dd x}
\newcommand*{\dxi}{\dd\xi}

\graphicspath{{./figures/}}

\newcommand*{\mk}{\mkern -1mu}
\newcommand*{\Mk}{\mkern -2mu}
\newcommand*{\mK}{\mkern 1mu}
\newcommand*{\MK}{\mkern 2mu}

\hypersetup{urlcolor=purple, linkcolor=blue, citecolor=red}


\title[Cauchy problem for the Schr\"odinger--Maxwell system]{Cauchy problem for the nonlinear Schr\"odinger equation coupled with the Maxwell equation}
\alttitle{Probl\`eme de Cauchy pour une \'equation de Schr\"odinger nonlin\'eaire coupl\'ee avec les \'equations de Maxwell}

\author[\initial{M.} \lastname{Colin}]{\firstname{Mathieu} \lastname{Colin}}
\address{University of Bordeaux, CNRS, Bordeaux INP, IMB, UMR 5251, F-33400, Talence, France \\
INRIA, IMB, UMR 5251, F-33400, Talence, France}
\email{mathieu.colin@math.u-bordeaux.fr}

\author[\initial{T.} \lastname{Watanabe}]{\firstname{Tatsuya} \lastname{Watanabe}}
\address{Department of Mathematics, Faculty of Science, \\
Kyoto Sangyo University, Motoyama, Kamigamo, Kita-ku, Kyoto-City, 603-8555, Japan}
\email{tatsuw@cc.kyoto-su.ac.jp}

\begin{abstract}
In this paper, we study the nonlinear Schr\"odinger equation coupled with the Maxwell equation. Using energy methods, we obtain a local existence result for the Cauchy problem.
\end{abstract}

\begin{altabstract}
Dans cet article, nous nous int\'eressons au couplage entre une \'equation de Schr\"odinger nonlin\'eaire et les \'equations de Maxwell. En utilisant des m\'ethodes d'\'energie, nous montrons que le probl\`eme de Cauchy est localement bien pos\'e.
\end{altabstract}

\subjclass{35L45, 35Q60, 35L70}
\keywords{Schr\"odinger--Maxwell system, Cauchy problem, symmetric hyperbolic system, energy method}

%\makeatletter
%\newcommand*\allowciteinthanks{\gdef\@thefnmark{}}
%\makeatother

\thanks{%\allowciteinthanks
The authors are grateful to anonymous referee for valuable comments and suggestions. The authors also would like to thank Professor P.~Antonelli for letting us know the reference~\cite{AAM}. The second author is supported by JSPS Grant-in-Aid for Scientific Research (C) (No.15K04970, No.18K03383).}

\begin{document}

\maketitle

\section{Introduction}

In this paper, we consider the following nonlinear Schr\"odinger equation coupled with Maxwell equation stated in $\R_+\times \R^3$ :
\begin{align}
i \psi_t +\Delta \psi &= e \phi \psi +e^2 | \BA |^2 \psi +2ie \nabla \psi \cdot \BA +ie \psi \Div \BA - g(|\psi|^2)\psi, \label{eq:1.1}
\\
\BA_{tt}-\Delta \BA &= e \IM (\bar{\psi} \nabla \psi)-e^2|\psi|^2 \BA -
\nabla \phi_t - \nabla \Div \BA, \label{eq:1.2}
\\
-\Delta \phi &= \frac{e}{2} |\psi|^2 + \Div \BA_t, \label{eq:1.3}
\end{align}
where $\psi: \R_+\times \R^3 \to \Com$, $\BA: \R_+\times \R^3 \to \R^3$, $\phi: \R_+ \times \R^3 \to \R$, $e \in \R$ and $i$ denotes the unit complex number, that is, $i^2=-1$. In this setting, $\psi$ is an electrically charged field and $(\phi,\BA)$ represents a gauge potential of an electromagnetic field. System~\eqref{eq:1.1}--\eqref{eq:1.3} describes the interaction of this Schr\"odinger wave function $\psi$ with the Maxwell gauge potential. The constant $e$ represents the strength of the interaction. For more details and physical backgrounds, we refer to~\cite{F}.

Since we are interested in the Cauchy Problem, let us consider the following set of initial data:
\begin{equation}\label{initial}
\psi(0,x) =\psi_{(0)}(x), \;\; \BA (0,x)=\BA_{(0)}(x), \;\; \BA_t(0,x)=\BA_{(1)}(x),
\end{equation}
where the regularity of each functions is given in Theorem~\ref{th1}. It is known that System~\eqref{eq:1.1}--\eqref{eq:1.3} has a so-called \emph{gauge ambiguity}. Namely if $(\psi,\BA, \phi)$ is a solution of~\eqref{eq:1.1}--\eqref{eq:1.3}, then $(\exp(i e \chi)\psi,\BA+\nabla \chi, \phi-\chi_t)$ is also a solution of~\eqref{eq:1.1}--\eqref{eq:1.3} for any smooth function $\chi:\R_+ \times \R^3 \to \R$. To push out this ambiguity, we adopt in the sequel the Coulomb gauge:
\begin{equation}\label{coulomb}
\Div \BA=0,
\end{equation}
which is propagated by the set of Equations~\eqref{eq:1.1}--\eqref{eq:1.3}. Indeed, if initially
\[
\Div \BA(0,{\cdot\,})= \Div \BA_t (0,{\cdot\,})=0,
\]
then~\eqref{coulomb} holds for all $t>0$. (See e.g.~\cite{CW1} for the proof.) In this setting, the last Equation~\eqref{eq:1.3} can be solved explicitly and the solution is given by
\[
\phi=\frac{e}{2}(-\Delta)^{-1}|\psi|^2,
\]
which imposes that
\[
\phi(0,x) = \frac{e}{2}(-\Delta)^{-1}|\psi_{(0)}(x)|^2.
\]
From~\eqref{coulomb}, we also observe that~\eqref{eq:1.1} can be written as
\begin{equation}\label{MS}
i \psi_t + L_A \psi - V(x) \psi + g(|\psi|^2)\psi =0,
\end{equation}
where $V$ is the non-local potential: $V(x)=\frac{e^2}{2} (-\Delta)^{-1} |\psi|^2$ and $L_A$ is the \emph{magnetic Schr\"odinger operator} which is defined by $\BA=(A_1,A_2,A_3)$ and
\begin{equation}\label{mag}
L_A \psi := \sum_{j=1}^3 \left(\frac{\partial}{\partial x_j}-ieA_j(x) \right)^2 \psi = \Delta \psi - 2ie \nabla \psi \cdot \BA - e^2 |\BA|^2 \psi.
\end{equation}

In this context, the two conserved quantities of the Schr\"odinger--Maxwell system are the charge $Q$ and the energy $E$:
\begin{equation}\label{charge}
Q(\psi)=\int_{\R^3} |\psi|^2 \,\dx,
\end{equation}
\vspace{-2\belowdisplayskip}
\begin{multline}\label{energy}
E(\psi,\BA,\phi)= \frac{1}{2}\int_{\R^3}
\big(|\nabla \psi -ie\BA \psi|^2+|\nabla \BA|^2 +|\partial_t \BA|^2\big) \,\dx\\
+\frac{e^2}{4} \int_{\R^3} \phi |\psi|^2 \,\dx -\int_{\R^3}G(|\psi|^2) \,\dx,
\end{multline}
where $G(t)=\int_0^t g(s)\ds$. To prove that~\eqref{charge} is formally conserved, one has to multiply Equation~\eqref{eq:1.1} by $\overline{\psi}$, integrate over $\R^3$ and take the imaginary part of the resulting equation. In a similar way, the conservation of~\eqref{energy} can be proved by multiplying~\eqref{eq:1.1}--\eqref{eq:1.3} by $\partial_t \overline{\psi}$, $\partial_t \BA$ and $\partial_t \phi$ respectively. This conserved quantities play a fundamental role if one wants to investigate the stability properties of such system, which is one of our main motivations. Indeed, in a previous paper~\cite{CW2}, we have showed that for small $e>0$, System~\eqref{eq:1.1}--\eqref{eq:1.3} admits a unique orbitally stable ground state of the form:
\begin{equation}\label{stand}
(\psi_{e,\omega}, \BA_{e,\omega}, \phi_{e,\omega}) := \left(\exp(i\omega t)u_{e,\omega}, \, \mathbf{0}, \,
\frac{e}{2} (-\Delta)^{-1} | u_{e,\omega}|^2 \right).
\end{equation}
In order to investigate the stability of such standing waves $(\psi_{e,\omega}, \BA_{e,\omega}, \phi_{e,\omega})$ as performed in~\cite{CW2}, it is necessary to prove that the Cauchy Problem~\eqref{eq:1.1}--\eqref{eq:1.3} is almost locally well-posed around $(\psi_{e,\omega}, \BA_{e,\omega}, \phi_{e,\omega})$.

In a previous paper~\cite{CW1}, we have proved the local existence of solutions for the nonlinear Klein--Gordon--Maxwell system in Sobolev spaces of high regularity. The method was to convert the Klein--Gordon--Maxwell system into a symmetric hyperbolic system and apply the standard energy estimate. Although our Schr\"odinger--Maxwell system~\eqref{eq:1.1}--\eqref{eq:1.3} looks similar, especially Equation~\eqref{eq:1.2} is completely the same, the usual reduction tools does not lead us to a symmetric hyperbolic system, which causes the necessity of a new strategy.

Let us also introduce results concerning the solvability of the Cauchy problem related to~\eqref{eq:1.1}--\eqref{eq:1.3}. In~\cite{BT}, \cite{NW}, the linear Schr\"odinger equation $(g \equiv 0)$ coupled with the Maxwell equations has been studied. Using the Strichartz estimate, the authors obtained the global well-posedness in the energy space. Recently in~\cite{AAM}, it was shown, by using the Strichartz estimate obtained in~\cite{NW}, that the system~\eqref{eq:1.1}--\eqref{eq:1.3} is locally well-posed in $H^2 \times H^{\frac{3}{2}} \times H^{\frac{1}{2}}$ and the global existence holds for finite energy weak solutions, when the nonlinear term $g$ is \emph{defocusing} (namely the case with $+|\psi|^{p-1}\psi$ in~\eqref{eq:1.1}). We also mention the paper~\cite{NT}, where the Cauchy problem of the Schr\"odinger--Maxwell system in the Lorentz gauge has been studied by using the energy method. On the other hand, a huge attention has been paid in the magnetic Schr\"odinger equation~\eqref{MS}. Especially in~\cite{M}, the local well-posedness for~\eqref{MS} in the energy space has been established in the case $V \equiv 0$. However, in this situation, the magnetic potential $\BA$ is \emph{given} and was assumed to be $C^{\infty}$, which cannot be expected a priori in our case. We also refer to~\cite{DFVV} for the Strichartz estimate for the magnetic Schr\"odinger operator~\eqref{mag} in the case $\BA \in L^2_{loc}(\R^3)$.

We mention that if we look for the standing wave~\eqref{stand}, we are led to the following non-local elliptic problem:
\begin{equation}\label{SP}
-\Delta u + \omega u + \left(\frac{e^2}{8 \pi|x|} * |u|^2 \right) u = g(|u|^2)u \quad \hbox{in} \ \R^3,
\end{equation}
which is referred as the \emph{Schr\"odinger--Poisson(--Slater)} equation. The existence of ground states of~\eqref{SP} as well as their orbital stability have been widely studied (see~\cite{AP}, \cite{BF}, \cite{BS}, \cite{CDSS}, \cite{K} and references therein). Finally, the orbital stability of standing waves for the magnetic Schr\"odinger equation~\eqref{MS} has been considered in~\cite{CE}, \cite{G}. Our study on the solvability of the Cauchy problem for~\eqref{eq:1.1}--\eqref{eq:1.3} and the result established in~\cite{CW2} enable us to generalize these previous results to the full Schr\"odinger--Maxwell system.

Before stating the main result of this paper, we introduce the following notations. As usual, $L^p(\R^3)$ denotes the usual Lebesgue space:
\[
L^p(\R^3)=\left\{ u \in \mathcal{S}'(\R^3) \ ; \ \| u \| _{L^p} <+\infty \right\},
\]
where
\[
\| u \|_{L^p}=\left(\int_{\R^3} |u(x)|^p \,\dx \right)^{\frac{1}{p}} \
\text{ if } 1\leq p<+\infty
\]
and
\[
\| u \|_{L^{\infty}} =\text{ess\,sup}\left\{ |u(x)| \ ; \ x\in\R^3 \right\}.
\]
We define the Sobolev space $H^s(\R^3)$ as follows:
\[
H^s(\R^3)=\left \{u \in \mathcal{S}'(\R^3) \ ; \ \| u \|_{H^s(\R^3)}^2 = \int_{\R^3} (1+|\xi| ^2)^s |\mathcal{F}(u)(\xi)|^ 2 \,\dxi < +\infty \right \},
\]
where $\mathcal{F}(u)(\xi)$ is the Fourier transform of $u$. We also introduce the homogeneous Sobolev space $\dot{H}^1(\R^3)$ as being the completion of $C^{\infty}_0(\R^3,\Com)$ for the norm $u \to \| |\xi| \mathcal{F}(u)(\xi) \|_{L^2(\R^3)}$. Recall that the space $\dot{H}^1(\R^3)$ is continuously embedded into $L^6(\R^3)$. Finally let $C(I,E)$ be the space of continuous functions from an interval $I$ of $\R$ to a Banach space $E$. For $1\leq j \leq 3$, we set $\partial_{x_j}=\frac{\partial}{\partial x_j}$ and $\partial_{t}=\frac{\partial}{\partial t}$. For $k\in \N^3$, $k=(k_1,k_2,k_3)$, we denote $D^k u=\partial_{x_1}^{k_1} \, \partial_{x_2}^{k_2} \, \partial_{x_3}^{k_3} u$ and for a non-negative integer $s$, $D^s$ denotes the set of all partial space derivatives of order $s$. Different positive constants might be denoted by the same letter $C$. We also denote by $\text{Re}(u)$ and $\text{Im}(u)$ the real part and the imaginary part of $u$ respectively.

We assume that $g$ satisfies
\begin{equation}\label{nonlinearity}
g \in C^{m+1}(\R,\R) \ \text{ and } g(0)=0,
\end{equation}
for some $m \in \N$ with $m \ge 2$, so that the function $W:\Com \to \Com$ defined by $W(u):=g(|u|^2)u$ satisfies $W\in C^{m+1}(\Com,\Com)$, $W(0)=W'(0)=0$. Some typical examples of the nonlinear term $g$ are the power nonlinearity $g(s)=\pm s^{\frac{p-1}{2}}$ with $[p]\ge 2m+3$ ($[p]$ denotes the integer part of $p$), or the \emph{cubic-quintic} nonlinearity $g(s)= s-\lambda s^2$ for $\lambda>0$, which frequently appears in the study of solitons in physical literatures. (See~\cite{RV} for example.) In this setting, we prove the following result.

\begin{theo}\label{th1}
Let $s$ be any integer larger than $\frac{3}{2}$ and assume that $\psi_0\in H^{s+2}(\R^3,\Com)$, $\BA_{(0)}\in H^{s+2}(\R^3,\R^3)$, $\BA_{(1)} \in H^{s+1}(\R^3,\R^3)$ with $\Div \BA_{(0)}=0$, $\Div \BA_{(1)}=0$ and $g$ satisfies~\eqref{nonlinearity}. Then there exist $T^*>0$ and a unique solution $(\psi,\BA,\phi)$ to System~\eqref{eq:1.1}--\eqref{eq:1.3} satisfying the initial condition~\eqref{initial} such that
\[
\psi \in C([0,T^*];H^{s+2}(\R^3)) \cap C^1([0,T^*]; H^s(\R^3)),
\]
\[
\BA \in C([0,T^*];H^{s+2}(\R^3)) \cap C^1([0,T^*];H^{s+1}(\R^3)),
\]
\[
\phi \in C([0,T^*];\dot{H}^1(\R^3)\cap L^{\infty}(\R^3)), \
\nabla \phi \in C([0,T^*];H^{s+1}(\R^3)),
\]
\[
\phi_t \in C([0,T^*];\dot{H}^1(\R^3) \cap L^{\infty}(\R^3)), \
\nabla \phi_t \in C([0,T^*];H^{s+1}(\R^3)).
\]
\end{theo}

The proof of Theorem~\ref{th1} is based on energy estimates and particularly, on the strategies developed in~\cite{Co1} and~\cite{CC}. Note also that to overcome the loss of derivatives embedded in Equation~\eqref{eq:1.2}, we use the original idea of Ozawa and Tsutsumi presented in~\cite{OT}.

The paper is organized as follows. In Section~\ref{sec2}, we transform System~\eqref{eq:1.1}--\eqref{eq:1.3} into a system to which we can apply the usual energy method. Section~\ref{sec3} is devoted to the proof of Theorem~\ref{th1}.



\section{Transformation of the equations}\label{sec2}

In this section, we transform the original System~\eqref{eq:1.1}--\eqref{eq:1.3} into a new symmetric system to which we can apply an energy method. In order to overcome the loss of derivatives contained in Equations~\eqref{eq:1.1}--\eqref{eq:1.2}, we introduce the following new unknowns (see~\cite{OT}):
\[
\Psi=\partial_t \psi \quad \hbox{and} \quad \Phi=\partial_t \phi.
\]

Let us first derive equations for $\Psi$ and $\Phi$. Differentiating Equation~\eqref{eq:1.1} with respect to $t$, one obtains
\begin{multline*}
i\Psi_t +\Delta \Psi = e \Phi \psi +e \phi \Psi +e^2 |A|^2 \Psi+ 2e^2 \psi \BA \cdot \BA_t +2i e \nabla \Psi \cdot \BA +2i e \nabla \psi \cdot \BA_t \\
-g'(|\psi|^2)(|\psi|^2\Psi+\psi^2\overline{\Psi})-g(|\psi|^2)\Psi.
\end{multline*}
Taking advantage of the new unknown $\Psi$, we also transform Equation~\eqref{eq:1.1} into an elliptic version
\begin{equation*}
i\Psi +\Delta \psi = e\phi \psi +e^2 |\BA|^2 \psi + 2 ie \nabla \psi\cdot \BA -g(|\psi|^2)\psi.
\end{equation*}
Moreover, we derive an equation for $\Phi$ by applying $\partial_t$ on Equation~\eqref{eq:1.3}:
\begin{equation*}
-\Delta \Phi =\frac{e}{2} \big(\psi \overline{\Psi}+\Psi \overline{\psi}\big).
\end{equation*}

Next, in order to ensure the Coulomb condition on $\BA$ for all $t>0$, we introduce the projection operator ${\mathbb P}$ on divergence free vector fields :
\begin{align*}
{\mathbb P} \;\; : \;\; \Big(L^2(&\R^3)\Big)^3 \longrightarrow \Big(L^2(\R^3)\Big)^3\\
&\BA\hspace{7mm}\longmapsto \hspace{4mm} {\mathbb P}\BA= \big(-\Delta\big)^{-1} \text{rot rot} \, \BA,
\end{align*}
so that if $\Div\BA=0$, then ${\mathbb P}\BA =\BA$. Thus applying ${\mathbb P}$ on Equation~\eqref{eq:1.2}, we derive
\begin{equation}\label{eq:2.1}
\BA_{tt}-\Delta \BA = {\mathbb P} \Big(e \IM (\bar{\psi} \nabla \psi)-e^2|\psi|^2 \BA -\nabla \Phi \Big).
\end{equation}
Note that any solution to~\eqref{eq:2.1} satisfying
\[
\Div\BA(0,{\cdot\,})=0 \quad \hbox{and} \quad
\Div\BA_t (0,{\cdot\,})=0,
\]
obviously satisfies
\[
\Div\BA(t,{\cdot\,})=0 \quad \hbox{for all} \ t>0.
\]
At this step, we have transformed System~\eqref{eq:1.1}--\eqref{eq:1.3} into
\begin{align}
i\Psi +\Delta \psi &= e\phi \psi +e^2 |\BA|^2 \psi + 2 ie \nabla \psi\cdot \BA -g(|\psi|^2)\psi, \label{eq:2.2}
\\
i\Psi_t +\Delta \Psi &= e \Phi \psi +e \phi \Psi +e^2 |\BA|^2 \Psi + 2e^2 \psi \BA \cdot \BA_t +2i e \nabla \Psi \cdot \BA \label{eq:2.3} \\
\nonumber& \qquad  +2i e \nabla \psi \cdot \BA_t -g'(|\psi|^2)(|\psi|^2\Psi+\psi^2\overline{\Psi}) - g(|\psi|^2)\Psi,
\\
\BA_{tt}-\Delta \BA &={\mathbb P} \Big(e \IM (\bar{\psi} \nabla \psi)-e^2|\psi|^2 \BA -\nabla \Phi \Big),\label{eq:2.4}
\\
-\Delta \phi &=\frac{e}{2} |\psi|^2, \label{eq:2.5}
\\
-\Delta \Phi &=\frac{e}{2} \big(\psi \overline{\Psi}+\Psi \overline{\psi}\big).\label{eq:2.6}
\end{align}
In order to take advantage of elliptic regularity properties, we transform Equations~\eqref{eq:2.2} by adding $-\alpha \psi$ ($\alpha>0$ will be chosen in Lemma~\ref{elliptic2} below) to both sides of the equation to obtain:
\begin{equation}\label{eq:2.7}
(-\Delta +\alpha)\psi = i \Psi -e\phi \psi -e^2 |\BA|^2 \psi - 2 ie \nabla \psi\cdot \BA +g(|\psi|^2)\psi+\alpha \psi.
\end{equation}
For simplicity, introduce $\mathcal{U}=(\phi,\Phi)$ and rewrite Equations~\eqref{eq:2.5} and~\eqref{eq:2.6} as
\begin{equation}\label{eq:2.8}
-\Delta \, \mathcal{U} = F_1(\psi,\Psi),
\end{equation}
where
\[
F_1(\psi,\Psi)=\frac{e}{2}
\binom{|\psi |^2}{\psi \overline{\Psi}+\Psi \overline{\psi}}.
\]
Equation~\eqref{eq:2.3} is then transformed into
\begin{equation}\label{eq:2.9}
i\partial_t \Psi +\Delta \Psi = 2i e \nabla \Psi \cdot \BA +2ie \nabla\psi \cdot \BA_t +F_2(\mathcal{U},\psi,\Psi,\BA,\BA_t),
\end{equation}
where
\begin{multline*}
F_2(\mathcal{U},\psi,\Psi,\BA,\BA_t) = e\Phi\psi +e\phi\Psi +e^2| \BA |^2\Psi +2e^2 \psi \BA\cdot \BA_t \\
-g'(|\psi|^2)(|\psi|^2\Psi+\psi^2\overline{\Psi})-g(|\psi|^2)\Psi.
\end{multline*}

It is then necessary to work with $\BA_t$ as new unknown. We recall first that $\BA=(a_1,a_2,a_3)$. To properly write the equations on $\BA$ and $\BA_t$, for $j=1,2,3$, $k=1,2,3$ and $\ell=1,2,3$, we introduce
\begin{align*}
p_{j,k}&= \partial_{x_k} a_j, \\
q_j&= \Delta a_j, \\
r_j&= \partial_t a_j, \\
\lambda_{j,k} &= \partial_{x_k} \Delta^{-1} \partial_t a_j, \\
\mu_{j,k,\ell} &= \partial_{x_\ell} \lambda_{j,k} = \partial_{x_\ell} \partial_{x_k} \Delta^{-1} \partial_t a_j, \\
\nu_{j,k}&= \Delta \lambda_{j,k}= \partial_{x_k} \partial_t a_j,\\
\tau_{j,k}&= \partial_t \lambda_{j,k}= \partial_{x_k} \Delta^{-1} \partial_t^2 a_j,
\end{align*}
and set $\mathcal{A}=(\mathcal{A}_1, \mathcal{A}_2, \mathcal{A}_3)$ with $\mathcal{A}_j: \R_+ \times \R^3 \to \R^{24}$ and
\[
\mathcal{A}_j={}^t
(a_j, p_{j,k}, q_j, r_j, \lambda_{j,k}, \mu_{j,k,\ell}, \nu_{j,k}, \tau_{j,k}).
\]
We also need to give some details on the projection operator ${\mathbb P}$. For that purpose, we introduce the Riesz transform ${\mathbb R}_j$ from $L^2(\R^3)$ to $L^2(\R^3)$ which is given by
\[
{\mathbb R}_j= \partial_{x_j} (-\Delta)^{-\frac{1}{2}} \quad \hbox{for} \ j=1,2,3.
\]
Then, ${\mathbb P}$ can be rewritten as ${\mathbb P}=({\mathbb P}_{j,m})_{1\leq j,m\leq 3}$ where
\[
{\mathbb P}_{j,m} =\delta_{j,m} +{\mathbb R}_j {\mathbb R}_m.
\]

Now we compute the equations for each components of $\mathcal{A}_j$. First by the definitions of $\mathcal{A}_j$, one finds that
\begin{align*}
\partial_t a_j &= \Delta \Delta^{-1} \partial_t a_j = \sum_{k=1}^3 \partial_{x_k} (\partial_{x_k} \Delta^{-1} \partial_t a_j) = \sum_{k=1}^3 \partial_{x_k} \lambda_{j,k}, \\
\partial_t p_{j,k} &= \partial_t \partial_{x_k} a_j = \Delta (\partial_{x_k} \Delta^{-1} \partial_t a_j) = \Delta \lambda_{j,k} = \sum_{\ell=1}^3 \partial_{x_\ell} (\partial_{x_\ell} \lambda_{j,k}) = \sum_{\ell=1}^3 \partial_{x_\ell} \mu_{j,k,\ell}, \\
\partial_t q_j &= \partial_t \Delta a_j = \sum_{k=1}^3 \partial_{x_k} \Delta (\partial_{x_k} \Delta^{-1} \partial_t a_j) =\sum_{k=1}^3 \partial_{x_k} \Delta \lambda_{j,k} =\sum_{k=1}^3 \partial_{x_k} \nu_{j,k}, \\
\partial_t r_j &= \partial_t^2 a_j = \Delta \Delta^{-1} \partial_t^2 a_j = \sum_{k=1}^3 \partial_{x_k} \partial_t (\partial_{x_k} \Delta^{-1} \partial_t a_j) \\
&= \sum_{k=1}^3 \partial_{x_k} \partial_t \lambda_{j,k} =\sum_{k=1}^3 \partial_{x_k} \tau_{j,k}.
\end{align*}
Next from Equation~\eqref{eq:2.4}, we have
\[
\partial_t^2 a_j = \Delta a_j + \sum_{m=1}^3 {\mathbb P}_{j,m}
\Big(e \IM (\bar{\psi} \partial_{x_m} \psi)-e^2|\psi|^2 a_m
-\partial_{x_m} \Phi \Big),
\]
which provides
\begin{gather*}
\begin{aligned}
\partial_t \lambda_{j,k} &= \partial_t \partial_{x_k} \Delta^{-1} \partial_t a_j = \partial_{x_k} \Delta^{-1} \partial_t^2 a_j \\
&= \partial_{x_k} \Delta^{-1} \Big(\Delta a_j + \sum_{m=1}^3 {\mathbb P}_{j,m}
\Big(e \IM (\bar{\psi} \partial_{x_m} \psi)-e^2|\psi|^2 a_m -\partial_{x_m} \Phi \Big) \Big) \\
&= \partial_{x_k} a_j + h_{j,k}^1 (\psi, \BA),
\end{aligned}\\[1.25em]
\begin{aligned}
\partial_t \mu_{j,k,\ell} &= \partial_t \partial_{x_\ell} \partial_{x_k} \Delta^{-1} \partial_t a_j = \partial_{x_\ell} \partial_{x_k} \Delta^{-1} (\partial_t^2 a_j) \\
&= \partial_{x_\ell} \partial_{x_k} \Delta^{-1}
\Big(\Delta a_j + \sum_{m=1}^3 {\mathbb P}_{j,m}
\Big(e \IM (\bar{\psi} \partial_{x_m} \psi)-e^2|\psi|^2 a_m -\partial_{x_m} \Phi \Big) \Big) \\
&= \partial_{x_\ell} p_{j,k} +h_{j,k,\ell}^2 (\psi,\BA),
\end{aligned}\\[1.25em]
\begin{aligned}
\partial_t \nu_{j,k} &= \partial_{x_k} \partial_t^2 a_j \\
&= \partial_{x_k}
\Big(\Delta a_j + \sum_{m=1}^3 {\mathbb P}_{j,m}
\Big(e \IM (\bar{\psi} \partial_{x_m} \psi)-e^2|\psi|^2 a_m -\partial_{x_m} \Phi \Big) \Big) \\
&= \partial_{x_k} q_j +h_{j,k}^3(\psi,\BA),
\end{aligned}
\end{gather*}
where $h^1_{j,k}$, $h_{j,k,\ell}^2$, $h^3_{j,k}$ are non-local functions defined as follows:
\begin{gather*}
h^1_{j,k}(\psi,\BA) =\partial_{x_k} \Delta^{-1}\sum_{m=1}^3 {\mathbb P}_{j,m}
\big(e\text{Im}(\overline{\psi}\partial_{x_m} \psi)-e^2|\psi|^2a_m -\partial_{x_m} \Phi \big),
\\
h^2_{j,k,\ell}(\psi,\BA) =\partial_{x_\ell} \partial_{x_k} \Delta^{-1}\sum_{m=1}^3{\mathbb P}_{j,m}
\big(e\text{Im}(\overline{\psi}\partial_{x_m} \psi)-e^2|\psi|^2a_m -\partial_{x_m} \Phi \big),
\\
h^3_{j,k}(\psi,\BA) =\partial_{x_k} \sum_{m=1}^3{\mathbb P}_{j,m}
\big(e\text{Im}(\overline{\psi}\partial_{x_m} \psi)-e^2|\psi|^2a_m -\partial_{x_m} \Phi \big).
\end{gather*}
Finally one has
\begin{align*}
\partial_t \tau_{j,k} &= \partial_{x_k} \Delta^{-1} \partial_t (\partial_t^2 a_j) \\
&= \partial_{x_k} \Delta^{-1} \partial_t
\Big(\Delta a_j + \sum_{m=1}^3 {\mathbb P}_{j,m}
\Big(e \IM (\bar{\psi} \partial_{x_m} \psi)-e^2|\psi|^2 a_m -\partial_{x_m} \Phi \Big) \Big).
\end{align*}
Computing separately each term of the right-hand side of the previous equation, we obtain
\begin{align*}
\partial_t(\overline{\psi}\partial_{x_m} \psi) &=\overline{\Psi}\partial_{x_m}\psi + \overline{\psi}\partial_{x_m} \Psi, \\
\partial_t(|\psi|^2a_m)&= (\Psi \overline{\psi}+\psi \overline{\Psi}) a_m + |\psi|^2 r_m.
\end{align*}
Moreover from~\eqref{eq:2.3} and~\eqref{eq:2.6}, one finds that
\begin{align*}
\partial_t (\partial_{x_m} \Phi) &= \partial_t \Big(\frac{e}{2} \partial_{x_m} (-\Delta)^{-1} (\psi \overline{\Psi} + \Psi \overline{\psi}) \Big) \\
&= \frac{e}{2} \partial_{x_m} (-\Delta)^{-1} (2|\Psi|^2+\psi \partial_t \overline{\Psi} +\overline{\psi} \partial_t \Psi) \\
&= e \partial_{x_m} (-\Delta)^{-1} |\Psi|^2 +e \partial_{x_m} (-\Delta)^{-1} \IM (i \partial_t \Psi \overline{\psi}) \\
&= e \partial_{x_m} (-\Delta)^{-1} \Big\{ |\Psi|^2 + \IM \Big(-\overline{\psi} \Delta \Psi +e \Phi |\psi|^2 +e \phi \overline{\psi} \Psi +e^2 |\BA|^2 \overline{\psi} \Psi \\
&\hspace{12em} + 2e^2 |\psi|^2 \BA \cdot \BA_t +2i e \overline{\psi} \nabla \Psi \cdot \BA +2i e \overline{\psi} \nabla \psi \cdot \BA_t \\
&\hspace{12em} -g'(|\psi|^2)|\psi|^2(\overline{\psi}\Psi+\psi\overline{\Psi}) - g(|\psi|^2) \overline{\psi} \Psi \Big) \Big\},
\end{align*}
from which we conclude that
\[
\partial_t \tau_{j,k}
= \partial_{x_k} r_j + h_{j,k}^4(\psi,\Psi,\BA, \BR),
\]
where $\BR=(r_1,r_2,r_3)$ and
\begin{align*}
&h^4_{j,k}(\psi,\Psi,\BA, \BR) \\
&=\partial_{x_k} \Delta^{-1}\sum_{m=1}^3 {\mathbb P}_{j,m}\Big(e\text{Im}(\overline{\Psi}\partial_{x_m}\psi + \overline{\psi}\partial_{x_m} \Psi) -e^2(|\psi|^2 r_m+(\Psi \overline{\psi}+\psi \overline{\Psi}) a_m)\Big) \\
&\quad -\partial_{x_k} \Delta^{-1}\sum_{m=1}^3 {\mathbb P}_{j,m} \Big[e \partial_{x_m} (-\Delta)^{-1} \Big\{ |\Psi|^2 + \IM \Big(-\overline{\psi} \Delta \Psi +e \Phi |\psi|^2 +e \phi \overline{\psi} \Psi \\
&\hspace{10em} +e^2 |\BA|^2 \overline{\psi} \Psi+ 2e^2 |\psi|^2 \BA \cdot \BR +2i e \overline{\psi} \nabla \Psi \cdot \BA +2i e \overline{\psi} \nabla \psi \cdot \BR \\
&\hspace{10em} -g'(|\psi|^2)|\psi|^2(\overline{\psi}\Psi+\psi\overline{\Psi}) - g(|\psi|^2) \overline{\psi} \Psi \Big) \Big\} \Big].
\end{align*}
The equation on $\mathcal{A}_j$ can be written as a symmetric system of the form
\[
\partial_t \mathcal{A}_j+\mathcal{M}_j(\nabla)\mathcal{A}_j +H_j(\psi,\Psi,\BA,\BR)=0 \quad (j=1,2,3),
\]
where $H_j={}^t (0,0,0,0,h_{j,k}^1, h_{j,k,\ell}^2, h_{j,k}^3, h_{j,k}^4)$, $\mathcal{M}_j(\nabla)= \sum_{k=1}^3 \widetilde{M_j} \partial_{x_k}$ are $24 \times 24$ symmetric matrices. Recalling that $\mathcal{A}_j=(a_j,p_j,q_j,r_j,\lambda_j,\mu_j,\nu_j,\tau_j)$, where $a_j$, $q_j$, $r_j$ are scalar functions, $p_j$, $\lambda_j$, $\nu_j$ and $\tau_j$ are functions with values in $\R^3$ and $\mu_j$ is a function with values in $\R^9$, $\mathcal{M}_j$ can be simply written by blocks:
\[
\mathcal{M}_j(\nabla)=
\left(
\begin{array}{cccccccc}
0&0&0&0&\nabla \cdot&0&0&0\\
0&0&0&0&0&\nabla \cdot&0&0\\
0&0&0&0&0&0&\nabla \cdot&0\\
0&0&0&0&0&0&0&\nabla \cdot \\
\nabla&0&0&0&0&0&0&0\\
0&N(\nabla)&0&0&0&0&0&0\\
0&0&\nabla&0&0&0&0&0\\
0&0&0&\nabla&0&0&0&0\\
\end{array}
\right),
\]
with
\[
N(\nabla)=
\left(
\begin{array}{ccc}
\nabla&0&0\\
0&\nabla&0\\
0&0&\nabla
\end{array}
\right).
\]
Note that $\widetilde{M_j}$ are $24 \times 24$ symmetric matrices whose components are all constants.



Thus from~\eqref{eq:2.7}, \eqref{eq:2.8} and~\eqref{eq:2.9}, we have transformed Equations~\eqref{eq:1.1}--\eqref{eq:1.3} into the following system:
\begin{align}
-\Delta \mathcal{U} &= F_1(\psi,\Psi), \label{eq:2.10}
\\
(-\Delta+\alpha)\psi + 2 ie \nabla \psi\cdot \BA &= i \Psi-e\phi \psi -e^2 |\BA|^2 \psi +g(|\psi|^2)\psi+\alpha \psi, \label{eq:2.11}
\\
i\partial_t \Psi +\Delta \Psi -2i e \nabla \Psi \cdot \BA &= 2ie \nabla\psi \cdot \BR +F_2(\mathcal{U},\psi,\Psi,\BA,\BR), \label{eq:2.12}
\\
0&=\partial_t \mathcal{A}_j+\mathcal{M}_j(\nabla)\mathcal{A}_j +H_j(\psi,\Psi,\BA,\BR). \label{eq:2.13}
\end{align}


\section{Solvability of the Cauchy Problem}\label{sec3}

The aim of this section is to prove Theorem~\ref{th1}. To this end, we use a fix-point argument on a suitable version of System~\eqref{eq:2.10}--\eqref{eq:2.13}. In this procedure, the necessary estimates follow from the application of the usual energy methods.

For $s \in \N$ with $s > \frac{3}{2}$, take an initial data
\[
\psi_{(0)} \in H^{s+2}(\R^3,\Com),
\]
\[
\BA_{(0)}=(a_{1(0)}, a_{2(0)}, a_{3(0)}) \in H^{s+2}(\R^3,\R^3),
\]
and
\[
\BA_{(1)}=(r_{1(0)},r_{2(0)},r_{3(0)}) \in H^{s+1}(\R^3,\R^3),
\]
satisfying
\[
\Div \BA_{(0)}=0, \;\Div \BA_{(1)}=0.
\]
Let us define $\Psi_{(0)} \in H^s(\R^3,\Com)$ by
\begin{equation}\label{Psi0}
\begin{gathered}
\Psi_{(0)}= i\big(\Delta \psi_{(0)} -e\phi_{(0)}\psi_{(0)}-e^2|\BA_{(0)}|^2 \psi_{(0)} -2ie\nabla \psi_{(0)}\cdot \BA_{(0)}-g(|\psi_{(0)}|^2)\psi_{(0)}\big),
\\
\Phi_{(0)}=\frac{e}{2}(-\Delta)^{-1} (\psi_{(0)} \overline{\Psi}_{(0)}
+\Psi_{(0)} \overline{\psi}_{(0)}),
\end{gathered}
\end{equation}
where $\phi_{(0)}= \frac{e}{2}(-\Delta)^{-1} | \psi_{(0)}|^2$.
We also put
\[
\mathcal{A}_{(0)}=(\mathcal{A}_{1(0)}, \mathcal{A}_{2(0)}, \mathcal{A}_{3(0)}) \in H^s(\R^3),
\]
for $i,j,k,l=1,2,3$,
\begin{gather*}
\mathcal{A}_{j(0)}= {}^t (a_{j(0)}, p_{j,k(0)}, q_{j(0)}, r_{j(0)},
\lambda_{j,k(0)}, \mu_{j,k,\ell(0)}, \nu_{j,k(0)}, \tau_{j,k(0)}),
\\
p_{j,k(0)}=\partial_{x_k} a_{j(0)}, \ q_{j(0)}=\Delta a_{j(0)}, \
\lambda_{j,k(0)}=\partial_{x_k} \Delta^{-1} r_{j(0)},
\\
\mu_{j,k,\ell(0)}= \partial_{x_k} \partial_{x_{\ell}} \Delta^{-1} r_{j(0)}, \
\nu_{j,k(0)}=\partial_{x_k} r_{j(0)}
\end{gather*}
and
\[
\tau_{j,k(0)}= \partial_{x_k} a_{j(0)} +\partial_{x_k} \Delta^{-1} \sum_{m=1}^3 \mathbb{P}_{j,m}
\big(e \IM (\overline{\psi}_{(0)}\partial_{x_m} \psi_{(0)}) -e^2 | \psi_{(0)}|^2 a_{m(0)}
- \partial_{x_m} \Phi_{(0)} \big).
\]
We introduce $R=2\big(\| \psi_{(0)} \|_{H^s}+\| \Psi_{(0)} \|_{H^s}+\| \mathcal{A}_{(0)} \|_{H^s}\big)$ and let $B(R)$ be the ball of radius $R$ in $C\big([0,T];(H^s(\R^3)^2\big)$ for $T>0$.

We prove the existence of a solution $(\mathcal{U}, \psi, \Psi, \mathcal{A}_j)$ of~\eqref{eq:2.10}--\eqref{eq:2.13} by the following procedure. Take $(\Psi, \mathcal{A}) \in B(R)$ with $\Div \BA=0$ arbitrarily and construct new functions $\mathcal{Q}$ and $\mathcal{B}=(\mathcal{B}_1,\mathcal{B}_2,\mathcal{B}_3)$ as follows.

First we define $\psi\in C([0,T];H^s(\R^3,\Com))$ by
\begin{equation}\label{psi}
\psi(t,x):=\psi_{(0)}(x)+ \int_0^t \Psi(s,x) \,\ds.
\end{equation}
Then by the construction of $\psi$, one finds that, for $T$ small enough,
\[
\| \psi\|_{L^{\infty}[0,T]; H^s)} \le R.
\]
Next let $\mathcal{U}\in C([0,T];\dot{H}^1(\R^3))$ be a solution to
\begin{equation}\label{eq:3.1}
-\Delta \mathcal{U} = F_1(\psi,\Psi).
\end{equation}
We note that $\mathcal{U} \in C([0,T]; L^{\infty}(\R^3))$ and $\nabla \mathcal{U} \in C([0,T];H^{s+1}(\R^3))$. (See Lemma~\ref{elliptic1} below.) Next we introduce the solution $\chi \in C([0,T];H^{s+2}(\R^3,\Com))$ of the following elliptic equation:
\begin{equation}\label{eq:3.2}
(-\Delta+\alpha) \chi + 2 ie \nabla \chi\cdot \BA = i\Psi -e\phi \psi -e^2 |\BA|^2 \psi +g(|\psi|^2)\psi+\alpha \psi.
\end{equation}
We now consider a linearized version of~\eqref{eq:2.12}--\eqref{eq:2.13}. We take $(\mathcal{Q},\mathcal{B}) \in H^s(\R^3) \times H^s(\R^3)$ solutions to
\begin{gather}\label{eq:3.3}
\begin{cases}
i\partial_t \mathcal{Q}+\Delta \mathcal{Q}- 2i e \nabla \mathcal{Q} \cdot \BA = 2ie \nabla\chi \cdot \BR+F_2(\mathcal{U},\chi,\Psi,\BA, \BR), \\
\mathcal{Q}(0,x)=\Psi_{(0)},
\end{cases}
\\ \label{eq:3.4}
\begin{cases}
\partial_t \mathcal{B}_j+\mathcal{M}_j(\nabla)\mathcal{B}_j+H_j(\chi,\Psi, \BA, \BR)=0,\\
\mathcal{B}_j(0,x)=\mathcal{A}_{j(0)}.
\end{cases}
\end{gather}
Let
\[
\mathcal{S}\;\; :\;\; (\Psi,\mathcal{A}) \longmapsto (\mathcal{Q},\mathcal{B}).
\]
Our strategy consists in showing that $\mathcal{S}$ is a contraction mapping on $B(R)$, provided that $T>0$ is sufficiently small and to prove that $\chi=\psi$, from which we obtain the existence of a solution $(\mathcal{U}, \psi, \Psi, \mathcal{A}_j)$ of~\eqref{eq:2.10}--\eqref{eq:2.13} and complete the proof of Theorem~\ref{th1}.

The proof is divided into 6 steps. We first recall the following classical lemma. (See e.g.~\cite[Proposition~2.1.1, p.~98]{AG} for the proof.)

\begin{lemm}\label{alhinac}
Let $u$, $v \in L^\infty(\R^3)\cap H^s(\R^3)$ for $s \in \N$. Then for all $(m_1,m_2)\in \N^3\times \N^3$ with $|m_1|+|m_2|=s$, one has
\[
\| D^{m_1}u D^{m_2}v \|_{L^2}
\leq C\big(\|u \|_{L^\infty} \| v \|_{H^s} +\| v \|_{L^\infty} \| u \|_{H^s}\big).
\]
\end{lemm}

\subsection*{Step 1: Solving the elliptic equation~(\ref{eq:3.1})}

\begin{lemm}\label{elliptic1}
There exists a unique solution $\mathcal{U}\in C([0,T];\dot{H}^1(\R^3))$ of~\eqref{eq:3.1}. Moreover, $\mathcal{U}=(\phi,\Phi)$ satisfies the following estimates.
\begin{align}
\| \nabla\phi \|_{L^\infty([0,T];H^{s+1})} \leq C_1(R), \;\;
\| \phi \|_{L^\infty([0,T];L^\infty)} \leq C_2(R), \label{est1}
\\
\| \nabla\Phi \|_{L^\infty([0,T];H^{s+1})} \leq C_3(R), \;\;
\| \Phi \|_{L^\infty([0,T];L^\infty)} \leq C_4(R), \label{est2}
\end{align}
where $C_1$, $C_2$, $C_3$ and $C_4$ are positive constants depending only on $R$.
\end{lemm}

\begin{proof}
First we note that the bilinear form
\[
a(u,v):=\int_{\R^3} \nabla u\cdot \nabla v \,\dx
\]
is continuous and elliptic on $\dot{H}^1(\R^3,\R) \times \dot{H}^1(\R^3,\R)$. Moreover since $\psi \in H^s(\R^3)$ and $\Psi \in H^s(\R^3)$, a direct computation gives
\[
\big\| |\psi|^2 \big\|_{L^{\frac{6}{5}}}
= \| \psi \|_{L^{\frac{12}{5}}}^2 \leq C \| \psi \|_{H^1}^2,
\]
\[
\| \psi \overline{\Psi}+ \Psi \overline{\psi} \|_{L^{\frac{6}{5}}}
\leq C \| \psi \|_{L^3} \| \Psi \|_{L^2} \leq C \| \psi \|_{H^1} \| \Psi \|_{L^2}.
\]
Then by the Sobolev embedding $L^{\frac{6}{5}}(\R^3) \hookrightarrow \big(\dot{H}^1(\R^3)\big)^*$ and the Lax--Milgram theorem, we deduce that there exists a unique solution $\mathcal{U} \in C([0,T];\dot{H}^1(\R^3))$ of~\eqref{eq:3.1}.

Next for $0 \le k \le s$, we apply $D^{k+1}$ to the first line of~\eqref{eq:3.1}, multiply the resulting equation by $D^{k+1} \phi$ and make an integration by parts to obtain
\begin{align*}
\| \nabla (D^{k+1} \phi) \|_{L^2}^2 &= \frac{e}{2} \left | \int_{\R^3} D^{k+1}|\psi |^2 D^{k+1} \phi \,\dx \right | \nonumber \\
& \le C \int_{\R^3} \big| D^{k}|\psi |^2 \big| |D^{k+2} \phi| \,\dx.
\end{align*}
Using the Leibniz rule, Lemma~\ref{alhinac} and the Schwarz inequality, one has
\begin{equation}\label{ellip}
\| \nabla \phi(t,{\cdot\,}) \|_{H^{k+1}}
\le C \| \psi (t,{\cdot\,}) \|_{H^k}^2 \quad
\hbox{for all} \ t \in [0,T].
\end{equation}
Summing up the inequalities~\eqref{ellip} from $k=0$ to $s$ and recalling the fact that $\| \psi \|_{L^\infty([0,T];H^s)}\leq R$, we obtain
\[
\| \nabla\phi \|_{L^\infty([0,T];H^{s+1})}\leq C_1(R),
\]
where $C_1(R)$ is a constant depending only on $R$.

Finally, the Sobolev embedding $W^{1,6}(\R^3) \hookrightarrow L^\infty(\R^3)$ provides that
\begin{align*}
\| \phi(t,{\cdot\,}) \|_{L^\infty} &\leq C \left(\sum_{k=1}^3 \| \partial_{x_k} \phi(t,{\cdot\,}) \|_{L^6} +\| \phi(t,{\cdot\,}) \|_{L^6} \right) \\
&\le C \left(\sum_{k=1}^3 \| \nabla (\partial_{x_k} \phi)(t,{\cdot\,}) \|_{L^2} +\| \nabla \phi(t,{\cdot\,}) \|_{L^2} \right),
\end{align*}
from which we deduce that there exists a constant $C_2(R)$ depending only on $R$ such that
\[
\| \phi \|_{L^\infty([0,T];L^\infty)} \leq C_2(R),
\]
which ends the proof of~\eqref{est1}. The proof of estimates~\eqref{est2} is similar and we omit the details.
\end{proof}

\subsection*{Step 2: Solving the elliptic equation~(\ref{eq:3.2})}

\begin{lemm}\label{elliptic2}
Suppose that $\BA \in H^s(\R^3,\R^3)$, $s > \frac{3}{2}$ and $\Div \BA=0$. Then for sufficiently large $\alpha>0$, the bilinear form
\[
b(u,v) := \int_{\R^3} (\nabla u \cdot \nabla \overline{v}
+\alpha u\overline{v}+ 2ie \nabla u\cdot \BA \overline{v}) \,\dx
\]
is hermitian, continuous and elliptic on $H^1(\R^3,\Com)\times H^1(\R^3,\Com)$.

As a consequence, there exists a unique solution $\chi(t,{\cdot\,}) \in H^1(\R^3,\Com)$ to~\eqref{eq:3.2} and there exists a constant $C_5(R)$ such that
\[
\|\chi \|_{L^\infty([0,T];H^{s+2})}\leq C_5(R).
\]
\end{lemm}

\begin{proof}
First we note that $b$ is hermitian by the condition $\Div \BA=0$. Indeed, one has
\begin{align*}
2ie \int_{\R^3} (\nabla u\cdot \BA) \overline{v} \,\dx &= -2ie \int_{\R^3} \Div \BA u\overline{v} \,\dx -2ie \int_{\R^3} (\nabla \overline{v} \cdot \BA) u \,\dx \\
&= \overline{2ie \int_{\R^3}(\nabla {v}\cdot \BA)\overline{u} \,\dx},
\end{align*}
from which it follows directly that $b(u,v)=\overline{b(v,u)}$. The continuity is a direct consequence of the Cauchy--Schwarz inequality and the fact that $\BA\in H^s(\R^3) \hookrightarrow L^\infty(\R^3)$. Finally for all $u\in H^1(\R^3,\Com)$, we have
\begin{align*}
\left | 2ie\int_{\R^3} (\nabla u\cdot \BA) \overline{u} \,\dx\right| &\leq 2e \| \BA \|_{L^\infty} \| \nabla u \|_{L^2} \| u \|_{L^2} \\
&\leq \frac{1}{2} \| \nabla u \|_{L^2}^2 +2e^2 \| \BA \|_{L^\infty}^2 \| u \|_{L^2}^2\\
&\leq \frac{1}{2} \| \nabla u \|_{L^2}^2 +Ce^2 \| \BA \|_{H^s}^2 \| u \|_{L^2}^2.
\end{align*}
Taking $\alpha \ge 2Ce^2 \| \BA \|_{H^s}^2$, one gets
\[
b(u,u) \geq \frac{1}{2} \| \nabla u \|_{L^2}^2 + \frac{\alpha}{2} \| u \|_{L^2}^2.
\]
This shows that $b$ is elliptic on $H^1(\R^3,\Com)\times H^1(\R^3,\Com)$.

Now since $\psi$, $\Psi\in H^s(\R^3)$ and $\phi\in \dot{H}^1(\R^3)\cap L^\infty(\R^3)$, it is obvious that $ i\Psi-e\phi \psi -e^2 |\BA|^2 \psi -g(|\psi|^2)\psi+\alpha \psi$ belongs to $L^2(\R^3)\hookrightarrow (H^1(\R^3))^*$. Then the Lax--Milgram theorem ensures the existence of a unique solution $\chi$ to~\eqref{eq:3.2} in $H^1(\R^3)$. Using the elliptic regularity theory and recalling that
\[
(\psi,\Psi) \in C([0,T],H^s(\R^3))^2,
\; \phi \in C(([0,T];\dot{H}^1(\R^3)\cap L^\infty(\R^3)),
\]
\[
\nabla \phi \in C([0,T];H^s(\R^3)),\; \BA \in C([0,T],H^s(\R^3)),
\]
one gets
\[
\| \chi \|_{L^\infty([0,T];H^{s+2})}\leq C_5(R),
\]
where $C_5(R)$ is a constant depending only on $R$. This ends the proof of\linebreak Lemma~\ref{elliptic2}.
\end{proof}

\subsection*{Step 3: Solving the Schr\"odinger equation~(\ref{eq:3.3})}

For convenience, we introduce the real form of Equation~\eqref{eq:3.3}. Denote $\mathcal{R}=(\mathcal{R}_1,\mathcal{R}_2) =(\RE \mathcal{Q}, \IM\mathcal{Q})$ and write
\begin{equation}\label{real}
\partial_t \mathcal{R} + \mathcal{J} \Delta \mathcal{R} - 2e \sum_{j=1}^3 \mathcal{K}_j (\BA) \partial_{x_j} \mathcal{R} =\mathcal{L}_1 (\nabla \chi, \BR) +\mathcal{L}_2 (\mathcal{U}, \chi, \Psi, \BA, \BR),
\end{equation}
\[
\mathcal{R}(0,x)= \big(\RE \Psi_0(x),\IM\Psi_0(x) \big),
\]
where
\begin{gather*}
\mathcal{J} =
\begin{pmatrix}
0&1\\
-1&0
\end{pmatrix},
\;\; \mathcal{K}_j(\BA) =
\begin{pmatrix}
{a}_j&0\\
0&{a}_j
\end{pmatrix},
\\
\mathcal{L}_1 (\nabla\chi, \BR)=
\begin{pmatrix}
\IM(2ie \nabla \chi \cdot \BR) \\
-\RE (2ie \nabla \chi \cdot \BR)
\end{pmatrix},
\\
\mathcal{L}_2 (\mathcal{U}, \chi, \Psi, \BA, \BR) =
\begin{pmatrix}
\IM F_2(\mathcal{U}, \chi, \Psi, \BA, \BR) \\
-\RE F_2(\mathcal{U}, \chi, \Psi, \BA, \BR)
\end{pmatrix}.
\end{gather*}
Now for $\varepsilon>0$, we consider a long-wave type regularization of~\eqref{real} (see~\cite{CG}):
\begin{equation}\label{regul}
\partial_t (1-\varepsilon \Delta) \mathcal{R}_\varepsilon +\mathcal{J} \Delta \mathcal{R}_\varepsilon -2e \sum_{j=1}^3 \mathcal{K}_j(\BA) \partial_{x_j} \mathcal{R}_\varepsilon =\mathcal{L}_1+\mathcal{L}_2,
\end{equation}
with $\mathcal{R}_\varepsilon(0) =(1-\varepsilon \Delta)^{-1}(\RE \Psi_0,\IM\Psi_0)$. Since Equation~\eqref{regul} is linear and contains differential operator in space of at most zero order, one can show that there exists a unique solution $\mathcal{R}_\varepsilon \in C([0,T];H^s(\R^3))$ to Equation~\eqref{regul}. Furthermore, we have the following estimate.

\begin{lemm}\label{regul1}
Let $\mathcal{R}_\varepsilon$ be the unique solution of Equation~\eqref{regul}. Then there exist constants $C_6(R)$, $C_7(R)$ independent of $\varepsilon$ such that
\[
\| \mathcal{R}_\varepsilon \|_{L^\infty([0,T];H^s)}
\leq e^{C_6(R)T} \| \Psi_0 \|_{H^s} +\left(e^{C_7(R)T}-1 \right)^{\frac{1}{2}}.
\]
\end{lemm}

\begin{proof}
We first begin with the $L^2$-estimate. We multiply~\eqref{regul} by $\mathcal{R}_\varepsilon$ and integrate over $\R^3$. Since $\mathcal{J}$ is skew-symmetric, one obtains
\begin{multline}\label{l20}
\frac{\partial}{\partial t} \left(\frac{1}{2} \int_{\R^3}
\big(| \mathcal{R}_\varepsilon|^2+\varepsilon|\nabla \mathcal{R}_\varepsilon|^2 \big)\,\dx \right)
\\
= 2e \int_{\R^3} \sum_{j=1}^3 \mathcal{K}_j(\BA) \partial_{x_j}
\mathcal{R}_\varepsilon \cdot \mathcal{R}_\varepsilon \,\dx +\int_{\R^3} \mathcal{L}_1 \cdot \mathcal{R}_\varepsilon \,\dx +\int_{\R^3} \mathcal{L}_2 \cdot \mathcal{R}_\varepsilon \,\dx.
\end{multline}
For $j=1,2,3$, we have from $\| \partial_{x_j} a_j \|_{H^s} \le R$ that
\begin{align}\label{l21}
\left| \int_{\R^3} \mathcal{K}_j(\BA) \partial_{x_j}
\mathcal{R}_\varepsilon \cdot \mathcal{R}_\varepsilon \,\dx \right| &=\left| \frac{1}{2}\int_{\R^3} a_j \partial_{x_j} |\mathcal{R}_\varepsilon |^2 \,\dx \right| =\left| \frac{1}{2}\int_{\R^3} \partial_{x_j}a_j |\mathcal{R}_\varepsilon |^2 \,\dx \right|
\nonumber \\
&\leq \frac{1}{2} \| \partial_{x_j}a_j \|_{L^\infty} \| \mathcal{R}_\varepsilon \|_{L^2}^2
\le C(R) \| \mathcal{R}_\varepsilon \|_{L^2}^2.
\end{align}
Since $\Psi$, $\BR \in H^s$ and $\BA \in H^{s+1}$, using Lemmas~\ref{elliptic1}-\ref{elliptic2}, one can also compute as follows :
\begin{align}\label{l22}
\left| \int_{\R^3} \mathcal{L}_1 (\nabla\chi, \BR) \cdot \mathcal{R}_\varepsilon \,\dx \right| & \leq C(R) \| \mathcal{R}_\varepsilon \|_{L^2}, \nonumber \\
\left| \int_{\R^3} \mathcal{L}_2 (\mathcal{U},\chi,\Psi,\BA, \BR)
\cdot \mathcal{R}_\varepsilon \,\dx \right| &\leq C(R) \| \mathcal{R}_\varepsilon \|_{L^2}.
\end{align}
Collecting~\eqref{l20}--\eqref{l22}, we derive
\[
\frac{\partial}{\partial t} \| \mathcal{R}_\varepsilon \|_{L^2}^2
\leq C(R) \| \mathcal{R}_\varepsilon \|_{L^2}^2 +C(R).
\]
By the Gronwall inequality and from
\[
\| \mathcal{R}_{\varepsilon}(0,{\cdot\,}) \|_{L^2}
= \| (1-\varepsilon \Delta)^{-1} (\RE \Psi_0,\IM\Psi_0) \|_{L^2}
\le \| \Psi_0 \|_{L^2},
\]
it follows that
\begin{align*}
\| \mathcal{R}_\varepsilon (t,{\cdot\,}) \|_{L^2}^2 & \le e^{C(R)t} \big(\| \mathcal{R}_\varepsilon (0,{\cdot\,}) \|_{L^2}^2 +1-e^{-C(R)t} \big) \\
&\le e^{C(R)T} \| \Psi_0 \|_{L^2}^2 + e^{C(R)T}-1 \\
& \le \left(e^{ \frac{C(R)T}{2}} \| \Psi_0\|_{L^2} +(e^{C(R)T}-1)^{\frac{1}{2}} \right)^2 \quad
\hbox{for all} \ t\in [0,T].
\end{align*}

Next we perform the $H^s$-estimate. We apply $D^s$ on~\eqref{regul}, multiply the resulting equation by $D^s \mathcal{R}_\varepsilon$, integrate over $\R^3$ and use the Gronwall inequality. We limit our attention to non-trivial terms. Recalling that $\chi \in C([0,T];H^{s+2})$ and using Lemma~\ref{elliptic2}, we obtain
\[
\left| \int_{\R^3} D^s \mathcal{L}_1(\nabla\chi, \BR)
\cdot D^s \mathcal{R}_\varepsilon \,\dx \right| \leq C(R) \| \mathcal{R}_\varepsilon \|_{H^s}.
\]
Moreover, one gets
\begin{align*}
\left| \int_{\R^3} \mathcal{K}_j(\BA) \partial_{x_j}D^s \mathcal{R}_\varepsilon
\cdot D^s \mathcal{R}_\varepsilon \,\dx \right| &= \left| \frac{1}{2} \int_{\R^3} a_j \partial_{x_j} |D^s\mathcal{R}_\varepsilon |^2 \,\dx \right | \\
&= \left| \frac{1}{2} \int_{\R^3} \partial_{x_j} a_j |D^s \mathcal{R}_\varepsilon |^2 \,\dx \right | \nonumber \\
&\leq \frac{1}{2} \| \partial_{x_j}a_j \|_{L^\infty} \| \mathcal{R}_\varepsilon \|_{H^s}^2.
\end{align*}
Arguing similarly as above, one finds that
\[
\| \mathcal{R}_\varepsilon \|_{L^\infty([0,T];H^s)}
\leq e^{C_6(R)T} \| \Psi_0 \|_{H^s} + (e^{C_7(R)T}-1)^{\frac{1}{2}},
\]
which ends the proof of Lemma~\ref{regul1}.
\end{proof}

Now we argue as in~\cite{BDS}, \cite{CG} and we perform the limit $\varepsilon \to 0$. By Lemma~\ref{regul1}, we know that $\mathcal{R}_{\varepsilon}$ is uniformly bounded in $L^{\infty}([0,T],H^s)$. From~\eqref{regul}, one also has
\begin{multline*}
\partial_t \mathcal{R}_{\varepsilon} = -(1-\varepsilon \Delta)^{-1} \mathcal{J} \Delta \mathcal{R}_{\varepsilon} +2e (1-\varepsilon \Delta)^{-1} \sum_{j=1}^3 \mathcal{K}_j(\BA) \partial_{x_j} \mathcal{R}_{\varepsilon} \\
 + (1-\varepsilon \Delta)^{-1} \mathcal{L}_1 + (1-\varepsilon \Delta)^{-1} \mathcal{L}_2.
\end{multline*}
This implies that
\[
\| \partial_t \mathcal{R}_{\varepsilon} \|_{H^{s-2}}
\le C(R) \| \mathcal{R}_{\varepsilon} \|_{H^s} +\| \mathcal{L}_1 \|_{H^{s-2}} + \| \mathcal{L}_2 \|_{H^{s-2}}
\le C \quad \hbox{for all} \ t \ge 0 \ \hbox{and} \ \varepsilon \in (0,1].
\]
Thus passing to a subsequence, we may assume that
\[
\mathcal{R}_{\varepsilon} \to \mathcal{R} \in L^{\infty}([0,T],H^s),
\partial_t \mathcal{R}_{\varepsilon} \to \partial_t \mathcal{R} \in L^{\infty}([0,T],H^{s-2})
\ \hbox{in the weak * topology}.
\]
From~\eqref{regul}, one can see that $\mathcal{R}$ is a solution of Equation~\eqref{real} and satisfies
\[
\| \mathcal{R} \|_{L^\infty([0,T];H^s)} \leq
e^{C_6(R)T} \| \Psi_0 \|_{H^s} + (e^{C_7(R)T}-1)^{\frac{1}{2}}.
\]
Moreover since $\mathcal{R}_{\varepsilon}(0) \to (\RE \Psi_0,\IM\Psi_0)$, we get $\mathcal{R}(0)=(\RE \Psi_0,\IM\Psi_0)$. We then deduce the existence of a solution $\mathcal{Q}$ to Equation~\eqref{eq:3.3} satisfying
\begin{equation}\label{esti1}
\| \mathcal{Q} \|_{L^\infty([0,T];H^s)} \leq e^{C_6(R)T} \| \Psi_0 \|_{H^s} + (e^{C_7(R)T}-1)^{\frac{1}{2}}.
\end{equation}

\subsection*{Step 4: Solving the symmetric system~(\ref{eq:3.4})}

First we note that it is straightforward to prove the existence of a unique solution $\mathcal{B}_j$ to Equation~\eqref{eq:3.4}. (We refer to~\cite[Proposition~1.2, p.~115]{AG} for the proof.) Furthermore, by using the Fourier transform $\mathcal{F}$, one has directly
\[
\mathbb{R}_j u =\mathcal{F}^{-1}\Big(i\frac{\xi_j}{|\xi |} \mathcal{F}(u)\Big)
\quad \hbox{for} \ j=1,2,3,
\]
from which we deduce that $\mathbb{R}_j$ and hence $\mathbb{P}_{j,m}$ are bounded from $L^2(\R^3)$ to $L^2(\R^3)$. As a consequence, using the fact $\Psi \in H^s$, $\BA \in H^{s+2}$, $\BR \in H^{s+1}$ and $\chi \in H^{s+2}$, one can prove that
\[
\| H_j(\chi, \Psi, \BA, \BR) \|_{L^{\infty}([0,T];H^s)} \le C(R).
\]
Thus applying the energy estimate to~\eqref{eq:3.4}, recalling that $M_j(\nabla)= \sum_{k=1}^3 \widetilde{M}_j \partial_{x_k}$ is symmetric and using the fact $\widetilde{M}_j$ consists of constant elements, we get
\[
\frac{\partial}{\partial t} \| \mathcal{B}_j \|_{H^s}^2
\le \| \mathcal{B}_j \|_{H^s}^2 +C(R).
\]
Then by the Gronwall inequality, we obtain the following estimate.

\begin{lemm}\label{hyper}
Let $\mathcal{B}_j$ be the unique solution of~\eqref{eq:3.4}. Then there exists a constant $C_8(R)>0$ such that
\begin{equation}\label{esti2}
\| \mathcal{B}_j \|_{L^{\infty}([0,T];H^s)}
\le e^{\frac{T}{2}} \| \mathcal{A}_{j(0)} \|_{H^s} +C_8(R) (e^T-1)^{\frac{1}{2}}.
\end{equation}
\end{lemm}

Collecting~\eqref{esti1} and~\eqref{esti2}, we can state the following result.

\begin{prop}\label{prop1}
There exists $\hat{T}>0$ such that for $0<T\leq \hat{T}$, $\mathcal{S}$ maps $B(R)$ into itself.
\end{prop}

\subsection*{Step 5: Contraction mapping}

Now we establish the following result.

\begin{prop}\label{prop2}
There exists $T^* \in (0,\hat{T}]$ such that $\mathcal{S}$ is a contraction mapping in the $L^\infty ([0,T^*];L^2(\R^3))$-norm.
\end{prop}

\begin{proof}
The proof is based on the fact that $s>\frac{3}{2}$ and on the fact that all the functions of Equations~\eqref{eq:3.1}--\eqref{eq:3.4} are Lipschitz with respect to their arguments. The proof is classical and we omit the details.
\end{proof}

From Propositions~\ref{prop1}, \ref{prop2} and by the contraction mapping principle, it follows that there exists a unique $(\Psi, \mathcal{A}) \in B(R)$ such that
\[
\mathcal{S}(\Psi,\mathcal{A})=(\Psi,\mathcal{A}),
\]
that is, $\Psi$ is the unique solution of the Schr\"odinger equation:
\[
\begin{cases}
i\partial_t {\Psi}+\Delta {\Psi}- 2i e \nabla {\Psi} \cdot \BA = 2ie \nabla\chi \cdot \BR +F_2(\mathcal{U},\chi,\Psi,\BA, \BR), \\
\Psi(0,x)=\Psi_{(0)},
\end{cases}
\]
and $\mathcal{A}_j$ $(j=1,2,3)$ is the unique solution to the symmetric system:
\[
\begin{cases}
\partial_t \mathcal{A}_j+\mathcal{M}_j(\nabla)\mathcal{A}_j +H_j(\chi, \Psi, \BA, \BR)=0, \\
\mathcal{A}_j(0,x)=\mathcal{A}_{j(0)}.
\end{cases}
\]
Since $\mathcal{A}(0,x)=\mathcal{A}_{(0)}$, we also have $\BA(0,x)=\BA_{(0)}$ and $\BA_t(0,x)=\BA_{(1)}$.

\subsection*{Step 6: Proof of Theorem~\ref{th1} completed}

Let us now go back to the original problem~\eqref{eq:1.1}--\eqref{eq:1.3}. To this end, we first remark that $\partial_t \psi=\Psi$ by~\eqref{psi}. Applying $\partial_t$ on Equation~\eqref{eq:3.2}, comparing the resulting equation with Equation~\eqref{eq:2.12} and recalling that $\BR= \partial_t \BA$, one obtains
\begin{equation}\label{eq:3.5}
-\Delta (\partial_t \chi-\partial_t \psi) + 2ie \nabla(\partial_t \chi-\partial_t \psi)\cdot \BA=0.
\end{equation}
By Lemma~\ref{elliptic2}, we know that the bilinear form
\[
{b}(u,v) =\int_{\R^3} (\nabla u\cdot \nabla \overline{v}+\alpha u\overline{v}
+2ie \nabla u\cdot \BA \overline{v}) \,\dx
\]
is hermitian, continuous and elliptic on $H^1(\R^3,\Com) \times H^1(\R^3,\Com)$, from which we deduce that Equation~\eqref{eq:3.5} has a unique solution. Since obviously 0 is a solution to Equation~\eqref{eq:3.5}, one has $\partial_t \chi=\partial_t \psi$. Moreover from~\eqref{psi}, it follows that $\psi(0,x)=\psi_{(0)}$ and hence $\phi(0,x)=\phi_{(0)}$ by the uniqueness of the solution of~\eqref{eq:3.1}. Thus substituting $t=0$ into~\eqref{eq:3.2}, we get
\[
-\Delta \chi(0) +2ie \nabla \chi(0) \cdot \BA_{(0)} =i \Psi_{(0)}-e\phi_{(0)} \psi_{(0)} -e^2 | \BA_{(0)}|^2 \psi_{(0)}
+g(| \psi_{(0)}|^2)\psi_{(0)}.
\]
By the definition of $\Psi_{(0)}$ given in~\eqref{Psi0}, one finds that
\[
-\Delta (\chi(0)-\psi_{(0)})
+2ie \nabla (\chi(0)-\psi_{(0)}) \cdot \BA_{(0)} =0,
\]
yielding that $\chi(0)=\psi_{(0)}$ by Lemma~\ref{elliptic2}.

Since $\partial_t\chi =\partial_t \psi$ and $\chi(0)=\psi_{(0)}$, it follows that $\chi=\psi$. As a consequence, $\psi$ is the unique solution to Equation~\eqref{eq:1.1}. Then from~\eqref{eq:3.2} and~\eqref{eq:3.3}, we conclude that $\BA$ and $\phi$ are the unique solutions to~\eqref{eq:1.2} and~\eqref{eq:1.3} respectively. Moreover by the uniqueness of~\eqref{eq:3.1}, one finds that $\partial_t \phi = \Psi$. Finally by Lemmas~\ref{elliptic1}, \ref{elliptic2} and from~\eqref{esti1}, \eqref{esti2}, $(\psi, \BA, \phi)$ has the desired regularity as stated in Theorem~\ref{th1}.



\bibliography{Colin}
\end{document}
